\section{Architecture}
\label{sec:architecture}
\label{sec:overview}
\label{sec:argument}

\label{sec:arch-verifiers}
\Ziren{} proves one statement, the execution of a MIPS32 image (\S\ref{sec:statement}), and outputs a \emph{compressed proof} of it. A \emph{native verifier} checks that proof as a program, natively or compiled to WebAssembly, as the consumers of the proofs published for real-time proving of Ethereum blocks do; it pays for the proof's size and its own running time. Algorithm~\ref{alg:architecture} states \Ziren{} up to that proof and Figure~\ref{fig:architecture} draws it. The subsections below annotate the algorithm and state its core as two protocols, the shard argument (Protocol~\ref{prot:shard}) and the block argument (Protocol~\ref{prot:block}), symbolic in the parameters of \S\ref{sec:argument-params}; Sections~\ref{sec:isa}--\ref{sec:recursion} construct their components and Section~\ref{sec:security} states what they establish.

Verifiers that count cost differently are served by \emph{applications} built on \Ziren{} (\S\ref{sec:applications}). Each verifies the compressed proof inside another proof system, starting from one shrink node: a wrap into a pairing-based SNARK for an \emph{on-chain verifier}, an Ethereum contract that pays gas, and a binary stage and narrow recursion for a \emph{garbled verifier}, a Boolean circuit that pays for its AND gates and input bits.

\begin{algorithm}[H]
\caption{\Ziren{}, up to the compressed proof. Lines~\ref{alg:arch-shards}--\ref{alg:arch-root} run as a pipeline (\S\ref{sec:arch-recursion}).}
\label{alg:architecture}
\DontPrintSemicolon
\KwIn{a program image $P$; input and hint streams $(\mathsf{in}, \mathsf{hint})$}
\KwOut{the compressed proof $\pi$}
$(\mathsf{pk}, \mathsf{vk}) \gets \mathsf{Setup}(P)$ \tcp*{commit the program and the preprocessed tables}
$(d_0, \dots, d_{N-1}) \gets \mathsf{Execute}(P, \mathsf{in}, \mathsf{hint})$ \tcp*{one compiled run, cut into shard descriptors}\label{alg:arch-exec}
\ForEach(\tcp*[f]{in parallel, one shard per accelerator}){$i \in \{0, \dots, N-1\}$}{\label{alg:arch-shards}
  $T_i \gets \mathsf{Replay}(d_i)$ \tcp*{re-execute shard $i$ to its event record}
  $\pi_i \gets \mathsf{ShardProve}(\mathsf{pk}, T_i)$ \tcp*{over $\KB$, committed under $\Poseidon$}\label{alg:arch-shard}
}
$\mathcal{S} \gets (\mathsf{Leaf}_{\sigma(\pi_i)}(\pi_i))_{i < N}$ \tcp*{one leaf key per quantised shard shape}\label{alg:arch-leaf}
\While{$|\mathcal{S}| > 1$}{
  $\mathcal{S} \gets \mathsf{Compose}(\mathcal{S})$ \tcp*{one level; a node takes $\leq 3$ contiguous children}
}
$\pi \gets \mathsf{Root}(\mathcal{S})$ \tcp*{arity one, under the compress schedule}\label{alg:arch-root}
\tcp{the native verifier runs $\mathsf{Verify}(\mathsf{vk}, h, \pi)$ (Protocol~\ref{prot:block})}
\end{algorithm}

\begin{figure}[t]
\centering
\resizebox{\textwidth}{!}{% The architecture: one proving stack, two verifiers.
\begin{tikzpicture}[
  font=\small,
  >={Latex[length=2mm]},
  box/.style={rectangle, draw=zirenblue, line width=0.6pt, minimum height=1.0cm, align=center, fill=white},
  stage/.style={box, text width=2.5cm},
  art/.style={rectangle, draw=zirenslate, line width=0.45pt, rounded corners=1.5pt, font=\scriptsize, minimum height=0.42cm, text width=2.75cm, inner xsep=1pt, align=center, fill=white},
  frame/.style={draw=zirenblue, line width=0.75pt, rounded corners=9pt},
  gframe/.style={frame, dashed},
  panel/.style={draw=zirenslate, dashed, line width=0.5pt},
  ver/.style={box, fill=zirenlight, line width=0.8pt, text width=2.6cm, minimum height=0.8cm},
  a/.style={->, line width=0.65pt, zirenorange},
  lab/.style={font=\footnotesize, zirenteal},
  tag/.style={font=\scriptsize\itshape, zirenteal},
  dot/.style={circle, draw=zirenblue, fill=white, line width=0.5pt, inner sep=0pt, minimum size=0.17cm},
]
% ---- program
\node[box, text width=2.3cm] (prog) at (0,2.2) {Guest program};
\node[box, inner sep=2pt, font=\fontsize{7}{8}\selectfont\ttfamily, align=left] (elf) at (0,0) {%
  \begin{tabular}{@{}l@{\hspace{0.6em}}l@{}}
  \textrm{\itshape clk} & \textrm{\itshape instruction}\\[-1pt]
  \hline
  $i$   & addu \$4,\$5,\$6\\[-2pt]
  $i{+}1$ & addiu \$5,\$4,0\\[-2pt]
  $i{+}2$ & lw \$6,4(\$29)\\[-2pt]
  $i{+}3$ & bne \$2,\$3,-4\\[-2pt]
  $i{+}4$ & syscall
  \end{tabular}};
\node[lab, below=0.08cm of elf] {MIPS32 image $P$};
\draw[a] (prog) -- node[lab, right] {compile} (elf);

% ---- execute and prove
\node[stage] (exec) at (3.9,0) {Execution\\{\scriptsize compiled run, cut into shard descriptors $d_i$}};
\foreach \k in {2,1} {
  \node[stage, minimum height=1.25cm] at ($(7.0,0)+(0.16*\k,0.16*\k)$) {};
}
\node[stage, minimum height=1.25cm] (shard) at (7.0,0) {ShardProve\\{\scriptsize replay $d_i$, prove; one shard per GPU}};
\node[stage, minimum height=1.25cm] (rec) at (10.3,0) {};
\node[font=\small, anchor=north] at (rec.north) {Recursion};
\begin{scope}[shift={(10.3,-0.3)}]
  \foreach \x [count=\j] in {-0.8,-0.4,0,0.4,0.8} \node[dot] (l\j) at (\x,-0.15) {};
  \node[dot] (c1) at (-0.4,0.18) {};
  \node[dot] (c2) at (0.6,0.18) {};
  \node[dot, fill=zirenblue] (rt) at (0.1,0.45) {};
  \foreach \j in {1,2,3} \draw[zirenblue, line width=0.4pt] (l\j) -- (c1);
  \foreach \j in {4,5} \draw[zirenblue, line width=0.4pt] (l\j) -- (c2);
  \draw[zirenblue, line width=0.4pt] (c1) -- (rt) (c2) -- (rt);
\end{scope}
\draw[a] (elf) -- node[lab, above, pos=0.3] {in, hint} (exec);
\draw[a] (exec) -- (shard);
\draw[a] ($(shard.east)+(0.16,0)$) -- node[lab, above] {$\pi_i$} (rec);

% artifacts panel
\node[art] (r1) at (4.1,2.55) {shard descriptors};
\node[art] (r2) at (7.0,2.55) {event records};
\node[art] (r3) at (9.9,2.55) {chips and buses};
\node[art] (r4) at (4.1,1.95) {\LogUp{}, zerocheck};
\node[art] (r5) at (7.0,1.95) {jagged \WHIR{}};
\node[art] (r6) at (9.9,1.95) {key allowlist};
\node[panel, fit=(r1)(r6), inner sep=4pt] (pan) {};
\node[tag, above=0.02cm of pan] {$\F = \KB$, $\Hash = \Poseidon$};
\node[gframe, solid, fit=(exec)(shard)(rec)(pan), inner xsep=8pt, inner ysep=14pt] (main) {};
\node[lab, font=\small, below=0.05cm of main] {execute and prove};

% ---- compressed proof and the native verifier
\node[box, text width=2.3cm] (cp) at (14.6,0) {Compressed\\proof $\pi$, $h$};
\node[ver] (nat) at (14.6,2.2) {Native verifier};
\draw[a] (main.east |- cp) -- (cp);
\draw[a] (cp) -- node[lab, right] {publish} (nat);

% ---- applications
\node[stage] (shrink) at (3.75,-5.0) {Shrink\\{\scriptsize verifies $\pi$}};
\node[tag, below=0.08cm of shrink, align=center] (tS) {$\F = \KB$\\$\Hash = \Poseidon$ or \Blake};
\node[stage] (wrap) at (7.25,-3.6) {Wrap\\{\scriptsize verifies the shrink}};
\node[stage] (snark) at (10.6,-3.6) {SNARK\\{\scriptsize Groth16 or PLONK over BN254}};
\node[tag, below=0.08cm of wrap] (tW) {\Poseidon{} over BN254};
\draw[a] (wrap) -- node[lab, above] {$\pi_{\mathrm w}$} (snark);
\node[stage, dashed] (bin) at (7.25,-6.4) {Binary stage\\{\scriptsize emulates \KB{} over bits}};
\node[stage, dashed] (nar) at (10.6,-6.4) {Narrow recursion\\{\scriptsize proves a recorded verifier tape}};
\node[tag, below=0.08cm of bin] (tB) {$\F = \mathbb F_{2^{128}}$, $\Hash = \Blake$};
\draw[a, dashed] (bin) -- node[lab, above] {$\pi_{\mathrm n}$} (nar);
\draw[a, dashed] (nar.south east) .. controls +(0.55,-0.85) and +(-0.55,-0.85) .. node[lab, below] (rp) {repeats} (nar.south west);
\draw[a] (shrink.east) -- ++(0.55,0) |- (wrap.west);
\draw[a, dashed] (shrink.east) ++(0.55,0) |- (bin.west);
\node[frame, fit=(shrink)(tS)(wrap)(snark)(tW)(bin)(nar)(tB)(rp)(main.west |- wrap)(main.east |- wrap), inner xsep=3pt, inner ysep=8pt] (apps) {};
\node[lab, font=\small, below=0.05cm of apps] {applications};
\node[ver] (oc) at (14.6,-3.6) {On-chain verifier};
\node[ver, dashed] (gv) at (14.6,-6.4) {Garbled verifier};
\draw[a] (snark.east) -- (snark.east -| apps.east) -- node[lab, above] {$\pi_{\mathrm{snark}}$} (oc);
\draw[a, dashed] (nar.east) -- (nar.east -| apps.east) -- node[lab, above] {$\pi_{\mathrm n}$} (gv);

% ---- the compressed proof feeds the shrink
\draw[a] (cp.south) -- ++(0,-1.75) -| node[lab, pos=0.25, above] {$\pi$} (shrink.north);
\end{tikzpicture}
}
\caption{One proving stack ends in a 274\,KiB compressed proof for a native verifier, and applications re-prove that proof for on-chain and garbled verifiers. Read left to right. \Ziren{} (top): a guest program compiles to a MIPS32 image; execution cuts one run into shards, which accelerators prove in parallel over \KB{} with \Poseidon{} commitments (the dashed panel lists the objects built along the way), and recursion composes the shard proofs into the compressed proof that a native verifier checks. Applications (bottom): a shrink node re-proves the compressed proof under the hash an application reads; the wrap closes it with a pairing-based SNARK for an on-chain verifier, and a binary stage over $\mathbb F_{2^{128}}$ and narrow recursion reach a garbled verifier.}
\label{fig:architecture}
\end{figure}

\subsection{Execution}
\label{sec:arch-execution}
\label{sec:executor}

Line~\ref{alg:arch-exec} runs once: a just-in-time compiler runs the guest and cuts the run into shards, each described by the values its reads observe and its boundary registers, and a worker replays its shard under the interpreter, which alone emits the event record the prover consumes, so no event record crosses a process boundary (\S\S\ref{sec:isa-record} and~\ref{sec:isa-emulator}). Nothing in the protocol depends on the compiler.

\subsection{The shard argument}
\label{sec:arch-shard}
\label{sec:argument-components}
\label{sec:argument-shard}

Line~\ref{alg:arch-shard} proves one shard. Chips, with one row per executed instruction, and buses arithmetise it (Section~\ref{sec:air}); their constraints and buses are designed to hold exactly when the shard is a correct piece of the execution, and the premises that remain open are listed in \S\ref{sec:security-relation}. A bus argument $\mathsf{BA}$ (\LogUp{}, \S\ref{sec:air-logup}) and a constraint argument $\mathsf{CA}$ (a zerocheck, \S\ref{sec:air-zerocheck}) reduce every claim to column evaluations at one point; a jagged commitment $\mathsf{JC}$ reduces those to one evaluation of one dense polynomial (\S\ref{sec:pcs}), which one \WHIR{} opening authenticates (\S\ref{sec:iopp}). Protocol~\ref{prot:shard} states the reduction. Each step ends with evaluation claims on committed columns, so the steps compose without the verifier reading a column, and the number, widths and heights of the chips are invisible above $\mathsf{JC}$.

\begin{protocol}[Shard argument]
\label{prot:shard}
\textbf{Statement:} a verifying key committing to the program and the
preprocessed tables, and the public values of one shard.\\
\textbf{Witness:} the shard's execution record.
\begin{enumerate}
  \item \textbf{Commit.} The prover arranges each chip's events into rows,
  places the columns of each commitment round into one dense vector and
  commits to it. Round order is fixed and the per-chip dimensions enter the
  transcript, so the geometry is public before any challenge is drawn
  (Section~\ref{sec:pcs}).
  \item \textbf{Buses.} The verifier draws the bus challenges. The prover runs
  the bus argument, which reduces every send and receive of every chip to
  column evaluations at one point (Section~\ref{sec:air-logup}).
  \item \textbf{Constraints.} The verifier draws the constraint challenges. The
  prover runs the constraint argument on the same point, so each column is
  opened once for both arguments (Section~\ref{sec:air-zerocheck}).
  \item \textbf{Jagged reduction.} The column claims of the two previous steps,
  together with explicit zero claims for the gaps in the committed geometry,
  are batched and reduced to one evaluation of the dense polynomial
  (Section~\ref{sec:pcs-jagged}).
  \item \textbf{Jagged assist.} The verifier does not materialise the indicator table
  the previous step reduced against. A second sumcheck over the public column
  boundaries certifies its evaluation, and any point-extension coordinates are sampled
  here, in an order that is itself part of the protocol
  (Section~\ref{sec:pcs-jagged-assist}).
  \item \textbf{Opening.} One batched opening authenticates that evaluation
  against every commitment of step~1 (Section~\ref{sec:iopp}).
\end{enumerate}
\textbf{Output:} a shard proof for the public values $\mathsf{pv}$ of
Definition~\ref{def:shard}.
\end{protocol}

\paragraph{Verifier checks.} No quantity a soundness bound depends on is read from the proof beyond the declared heights, which the configuration caps (a column is at most $2^{22}$ rows, \S\ref{sec:pcs-geometry}), and every algebraic identity the native verifier checks is also checked by the recursive verifier. The verifier (i) absorbs the verifying key, the public values and the per-chip dimensions before drawing any challenge; (ii) checks every sumcheck round of the bus argument, the constraint argument, the jagged sumcheck and the jagged assist, and each one's closing identity, evaluating the constraint polynomials itself; (iii) checks Equation~\eqref{eq:stacking-bind}; (iv) runs the \WHIR{} verifier with the query counts, folding factors and grinding bits of the schedule; and (v) checks that the public values are well formed. Preprocessed row counts come from the verifying key and the schedule from the verifier's configuration.

\paragraph{Costs.} With $A$ the committed area of a shard, $n_{\mathrm{int}}$ its number of bus interactions and $K$ its number of columns, the prover performs $O(A + n_{\mathrm{int}})$ field operations outside the commitment and $O(A/\rho)$ hashing and transform work inside it. The verifier checks $O(\log^2 n_{\mathrm{int}})$ sumcheck rounds for the layered lookup circuit and $O(\log A)$ for each of the other three sumchecks, evaluates each chip's constraints once, spends $O(mK)$ on the jagged assist's closing identity (Theorem~\ref{thm:jagged-assist}), and checks the \WHIR{} queries, whose count is fixed by the schedule. The schedule sets the proof size (\S\ref{sec:perf-proofsize}), and Section~\ref{sec:performance} gives the concrete values.

\subsection{The block argument}
\label{sec:arch-recursion}
\label{sec:argument-block}
\label{sec:perf-pipeline}

Lines~\ref{alg:arch-leaf}--\ref{alg:arch-root} compose the shard proofs into one proof of the arithmetised relation $\mathcal R_{\mathrm{arith}}$ of \S\ref{sec:security-relation}, which \S\ref{sec:verification} connects to the relation $\mathcal R$ of \S\ref{sec:statement}. Leaves and compose nodes are recursion programs proved with the same shard argument, and each checks its children's keys against an allowlist enumerated from the machine's declared shapes (\S\ref{sec:recursion}). Protocol~\ref{prot:block} states these lines, and the native verifier, with the conditions each node checks.

\begin{protocol}[Block proof, end to end]
\label{prot:block}
\textbf{Public parameters:} the field $\F$ and extension $\EF$, the $\Poseidon{}$ transcript and Merkle trees, the \WHIR{} schedule (Table~\ref{tab:whir-schedule}), the bucket rule (\S\ref{sec:pcs-geometry}), the compose arity (at most three), and the allowlist root $\mathsf{root}_K$ of the admissible leaf and compose keys (Algorithm~\ref{alg:enum}).

\smallskip
\noindent\textbf{Setup}$(P)$. Commit to the program ROM of image $P$ and to the preprocessed \texttt{Program}, \texttt{Byte} and \texttt{Range} tables; output $(\mathsf{pk}, \mathsf{vk})$.

\smallskip
\noindent\textbf{Prove}$(\mathsf{pk}, P, \mathit{in}, \mathit{hint})$.
\begin{enumerate}
  \item \textbf{Execute and shard.} Execute $P$ on $(\mathit{in}, \mathit{hint})$ and cut the run into shards $T_0, \dots, T_{N-1}$, shard $T_i$ having index $i+1$ and its record emitted by re-executing it (\S\ref{sec:isa-emulator}), at the first of the clock fence, the shard size, a chip's height budget and the area budget; each shard's public values $\mathsf{pv}_i$ are those of Definition~\ref{def:shard} (\S\ref{sec:isa-record}).
  \item \textbf{Prove the shards.} As each record is produced, prove it with the shard argument (Protocol~\ref{prot:shard}): commit, bus argument, zerocheck, jagged reduction and assist, and one \WHIR{} opening, giving $\pi_i$.
  \item \textbf{Leaves.} Quantise each shard's geometry into its class $\sigma_i = (\text{cluster}, \mathit{first}, b_{\mathrm p}, b_{\mathrm m}, \mathsf{pad}(b_{\mathrm m}))$. The leaf program of that class verifies $\pi_i$, or a few fused shard proofs, against $\mathsf{vk}$ and outputs the compose tuple of the range (\S\ref{sec:recursion}).
  \item \textbf{Compose.} As soon as contiguous ranges are complete, merge up to three of them with a compose node, which performs the checks of Algorithm~\ref{alg:compose}.
  \item \textbf{Root.} Once one node covers $T_0, \dots, T_{N-1}$, a compose node of arity one over it, proved under the compress schedule, performs the root checks of Algorithm~\ref{alg:compose} and sets $\mathsf{complete}$; the first shard's leaf has already required $\mathsf{entry}_0 = \mathsf{start}(\mathsf{vk})$ and starting cursors of zero (\S\ref{sec:recursion}). Its proof, with the recursion key $\mathsf{vk}^{\mathrm{rec}}$ that verifies it, is the compressed proof $\pi$.
\end{enumerate}

\noindent\textbf{Verify}$(\mathsf{vk}, h, \pi)$.
\begin{enumerate}
  \item Verify $\pi$ under $\mathsf{vk}^{\mathrm{rec}}$ with the shard verifier of \S\ref{sec:argument-shard} and the compress schedule, which the verifier fixes from the stage it expects and never reads from the proof.
  \item Require $\mathsf{vk}^{\mathrm{rec}}$ to be in the allowlist, the exported root to equal $\mathsf{root}_K$, and the completion flag to be set (\S\ref{sec:security-binding}).
  \item Require the exported guest key to equal $\mathsf{vk}$ and the exported output digest to equal $h$; accept.
\end{enumerate}
\end{protocol}

Steps~2 to~4 of \textbf{Prove} overlap: a shard is proved as soon as its record exists and a compose node as soon as its children are proved, with trace generation and proof-of-work on the accelerators. Only the root uses the schedule tuned for size (\S\ref{sec:perf-proofsize}). The start-up interval and the serial composes over the last shards remain (Table~\ref{tab:cards}).

\subsection{Parameters}
\label{sec:argument-params}

The protocol is symbolic in the following parameters; numeric values are fixed where each component is built.
\begin{itemize}
  \item The base field $\F$ and the challenge extension $\EF$, whose size bounds every term of the form $\deg/|\EF|$; the permutation of the transcript and commitment trees, with its round counts (\S\ref{sec:prelim}), and the trees' arity.
  \item The committed geometry: the stacking height $H=2^{\ell}$, the row cube $2^{n}$ and the bucket rule (\S\ref{sec:pcs-geometry}), which together fix the set of recursion programs a prover can produce.
  \item The proximity test: the initial rate $\rho$, per-round folding factors, query counts, out-of-domain samples, the grinding bits of each query phase and of the bus challenges (Table~\ref{tab:whir-schedule}), and the decoding regime.
  \item Recursion: the compose arity and the height of the tree of admissible keys (\S\ref{sec:recursion}).
  \item The machine: the shard size, clock fence and per-chip height budgets (\S\ref{sec:isa}), and the guest-accelerator set, which fixes the chip set.
  \item The cross-shard digest: the curve and its extension, the interaction bound $N_{\mathrm{int}}$ and the multiplicity bound $B$ (\S\ref{sec:air-digest-security}).
\end{itemize}

\subsection{Applications}
\label{sec:applications}

Figure~\ref{fig:architecture} (bottom) draws the two applications built on \Ziren{}; neither changes the statement, and each output attests the same execution as $\pi$.


\paragraph{The shrink.}\label{sec:arch-shrink} A shrink node verifies the compressed proof; it is proved under \Poseidon{} for the wrap and under \Blake{}~\citep{blake3}, commitments and transcript alike, for garbled verification, so garbled verification changes the hash before it changes the field (\S\ref{sec:recursion-shrink-wrap}).

\paragraph{On-chain verification.}\label{sec:arch-wrap} An on-chain verifier is an Ethereum contract; it pays gas, which precompiled BN254 pairings make small and independent of the statement only for a pairing-based argument over that curve with a constant-size proof~\citep{groth16,plonk}. A wrap node re-proves the shrink so that a Groth16 or PLONK circuit over BN254 verifies it while hashing in its native field, and the on-chain verifier checks that circuit's proof with a constant number of pairings (\S\ref{sec:recursion-shrink-wrap}).

\paragraph{Garbled verification.}\label{sec:arch-garbled} A garbled verifier is a Boolean circuit evaluated under garbling~\citep{yao86}, as in protocols that check a computation on Bitcoin only when it is disputed~\citep{bitvm}. Evaluating a garbled verifier on an invalid proof reveals a secret that serves as the fraud proof~\citep{glock}; later protocols build bridges on this~\citep{bitvm3}, make the verification of pairing-based SNARKs far cheaper~\citep{argo-mac,babe}, and shrink the on-chain footprint of malicious security~\citep{mosaic}. They garble the verifier of a pairing-based SNARK or of a designated-verifier variant of one, where this application garbles a hash-based verifier over a binary field. A garbled verifier pays for its AND gates and for its input bits, while XOR gates are free~\citep{free-xor,half-gates}. Over a 31-bit prime field an addition is a carry chain and a multiplication a reduction modulo $p$; over $\mathbb F_{2^{128}}$ an addition is a XOR, so this application moves the arithmetic to a binary field and the hash to one built from word additions, XORs and rotations. The binary stage commits bit columns as packed elements of $\mathbb F_{2^{128}}$ and opens them through ring switching~\citep{binius,fri-binius}, emulating the shrink verifier's \KB{} arithmetic over bits.

Its verifier is not rewritten as a circuit. A recorder runs it over a value type that appends every multiplication, inversion, bit decomposition and \Blake{} hash to a tape, and every comparison as the condition it found, so the tape is fixed by the proof's shape and key and enforces the control flow of the run that recorded it. A narrow machine native to $\mathbb F_{2^{128}}$ proves a run of the tape, with a ledger of read cells and tables for arithmetic, bit rewiring and \Blake{}, and narrow recursion repeats the step on each narrow proof's own verifier. The garbled circuit evaluates the last recorded tape on the bits of the last narrow proof; it pays AND gates for field products, \Blake{} compressions and the per-claim tensor algebra of ring switching, and takes as input every value the proof opens and every Merkle path. The narrow machine follows from that cost: few tables, narrow ones, none reading its next row, power-of-two widths, and a low-rate opening.
